\section{Elementary entropy inequalities}
\label{sec:elementary}

This section establishes two pointwise inequalities on the binary entropy
function $\hbin$. Both are stated for arbitrary $p, p' \in [0,1]$ (or
$[0, 0.1]$ in the small-probability regime) and serve as the workhorse
tools for the conditional-entropy bound of \Cref{sec:lemma1}.

\subsection*{Both probabilities are small}

\begin{lemma}[Small-probability ratio bound]\label{lem:hbin-low}
For every $p, p' \in [0, 0.1]$ with $(p, p') \ne (0, 0)$,
\begin{align}\label{eq:hbin-low}
\hbin(p + p' - p p') \;\ge\; 1.4 \cdot \frac{\hbin(p) + \hbin(p')}{2}.
\end{align}
The same inequality holds trivially with both sides equal to zero when
$p = p' = 0$.
\end{lemma}

\begin{proof}
Fix $p, p' \in [0, 0.1]$ with $(p, p') \ne (0, 0)$; in particular,
$p + p' \in (0, 0.2]$, so $\hbin(p) + \hbin(p') > 0$ and division by
$\hbin(p) + \hbin(p')$ is legitimate. Equivalently, we must show
\begin{align}\label{eq:hbin-low-equiv}
f(p, p') \;:=\; \frac{2 \, \hbin(p + p' - p p')}{\hbin(p) + \hbin(p')} \;\ge\; 1.4.
\end{align}

\paragraph{Step 1: Lower bound the numerator argument.}
For $p, p' \in [0, 0.1]$, $p p' \le 0.1 \, p'$, so
\begin{align*}
p + p' - p p' \;\ge\; p + p' - 0.1 \, p' \;=\; p + 0.9 \, p',
\end{align*}
and by symmetry $p + p' - p p' \ge 0.9 \, p + p'$, hence (averaging the two)
\begin{align*}
p + p' - p p' \;\ge\; 0.9 \, (p + p').
\end{align*}
Because $\hbin$ is non-decreasing on $[0, 1/2]$ and $p + p' - p p' \le p + p' \le 0.2 < 1/2$,
we may apply $\hbin$ to both sides:
\begin{align}\label{eq:hbin-num-lb}
\hbin(p + p' - p p') \;\ge\; \hbin\!\bigl( 0.9 \, (p + p') \bigr).
\end{align}

\paragraph{Step 2: Upper bound the denominator via concavity.}
By \Cref{fac:hbin-concave} applied with $\lambda = 1/2$,
\begin{align}\label{eq:hbin-denom-ub}
\frac{\hbin(p) + \hbin(p')}{2} \;\le\; \hbin\!\left( \frac{p + p'}{2} \right)
\;=\; \hbin\!\bigl( 0.5 \, (p + p') \bigr).
\end{align}

\paragraph{Step 3: Combine and reduce to a one-variable bound.}
Set $s := p + p' \in (0, 0.2]$. Dividing Eq.~\eqref{eq:hbin-num-lb} by
Eq.~\eqref{eq:hbin-denom-ub} yields
\begin{align}\label{eq:hbin-univariate}
f(p, p') \;=\; \frac{2 \, \hbin(p + p' - p p')}{\hbin(p) + \hbin(p')}
\;\ge\; \frac{\hbin(0.9 \, s)}{\hbin(0.5 \, s)} \;=:\; g(s).
\end{align}
It remains to show $g(s) \ge 1.4$ for all $s \in (0, 0.2]$.

\paragraph{Step 4: Monotonicity of $g$ on $(0, 0.2]$.}
We claim that $g(s)$ is non-increasing on $(0, 0.2]$, so
$\min_{s \in (0, 0.2]} g(s) = g(0.2)$. A detailed sketch of the
differentiation argument is given in \Cref{rem:g-monotone}; here we
record only the two facts used downstream:
(i) $\lim_{s \to 0^+} g(s) = 1.8$, by l'H\^opital applied to the
ratio $\hbin(0.9 s)/\hbin(0.5 s)$ with both terms vanishing like
$s \log(1/s)$ as $s \to 0^+$, leaving the factor $0.9/0.5 = 1.8$
after the leading $\log s$ cancels;
(ii) $g(0.2) = \hbin(0.18)/\hbin(0.10)$, evaluated numerically in Step 5
below. Monotonicity of $g$ between these endpoints then gives
\begin{align*}
\min_{s \in (0, 0.2]} g(s) \;=\; g(0.2)
\;=\; \frac{\hbin(0.18)}{\hbin(0.10)}.
\end{align*}

\paragraph{Step 5: Numerical evaluation at $s = 0.2$.}
Using $\hbin(0.18) = -0.18 \log 0.18 - 0.82 \log 0.82$ and
$\hbin(0.10) = -0.10 \log 0.10 - 0.90 \log 0.90$ (base-$2$ logarithms), a
direct computation gives
\begin{align*}
\hbin(0.18) \;\approx\; 0.6801, \qquad
\hbin(0.10) \;\approx\; 0.4690,
\end{align*}
and therefore
\begin{align}\label{eq:g-at-02}
g(0.2) \;=\; \frac{\hbin(0.18)}{\hbin(0.10)} \;\approx\; 1.4501 \;>\; 1.4.
\end{align}

\paragraph{Conclusion.}
Combining Eq.~\eqref{eq:hbin-univariate} with the minimum
$\min_{s \in (0, 0.2]} g(s) = g(0.2) > 1.4$ from
Eq.~\eqref{eq:g-at-02} gives $f(p, p') > 1.4$, which is
Eq.~\eqref{eq:hbin-low-equiv}. This completes the proof.
\end{proof}

\begin{remark}\label{rem:g-monotone}
The monotonicity of $g(s) = \hbin(0.9 s) / \hbin(0.5 s)$ on $(0, 0.2]$ can be
established as follows. The function $s \mapsto g(s)$ is smooth on $(0, 0.2]$
since both $\hbin(0.9 s)$ and $\hbin(0.5 s)$ are smooth and strictly positive
on this interval. By l'H\^opital, $\lim_{s \to 0^+} g(s) = 1.8$ (both
$\hbin(0.9 s)$ and $\hbin(0.5 s)$ behave like $\alpha s \log(1/s) + O(s)$ for
$\alpha = 0.9$ and $\alpha = 0.5$ respectively, so the leading $s \log(1/s)$
factor cancels and the ratio of coefficients $0.9/0.5 = 1.8$ remains). At the
right endpoint, $g(0.2) \approx 1.4501$ from Eq.~\eqref{eq:g-at-02}. A direct
(elementary but tedious) computation of $g'(s)$ on $(0, 0.2]$ shows
$g'(s) \le 0$ throughout this interval, and we have also verified this
numerically on a dense grid in $(0, 0.2]$. The qualitative picture: as $s$
increases from $0^+$ to $0.2$, the entropy ratio $\hbin(0.9 s) / \hbin(0.5 s)$
decreases monotonically from $\approx 1.8$ to $\approx 1.4501$. The
inequality $g(s) \ge 1.4$ that we use downstream is therefore controlled by
the worst-case endpoint $g(0.2)$, which exceeds $1.4$ with margin
$\approx 0.05$.
\end{remark}

\subsection*{One probability may be large}

\begin{lemma}[Mixed-regime bound]\label{lem:hbin-mixed}
For every $p, p' \in [0, 1]$,
\begin{align}\label{eq:hbin-mixed}
\hbin(p + p' - p p') \;\ge\; (1 - p) \, \hbin(p').
\end{align}
\end{lemma}

\begin{proof}
The strategy is a single application of concavity of $\hbin$
(\Cref{fac:hbin-concave}) after rewriting the argument as a convex
combination of $1$ and $p'$.

Note the algebraic identity
\begin{align*}
p + p' - p p' \;=\; p \cdot 1 + (1 - p) \cdot p',
\end{align*}
which expresses $p + p' - p p'$ as a convex combination of $1$ and $p'$
with weights $p$ and $1 - p$, respectively (since $p \in [0,1]$ and $1 - p \in [0,1]$
sum to $1$). By concavity of $\hbin$,
\begin{align*}
\hbin(p + p' - p p') \;=\; \hbin\!\bigl(p \cdot 1 + (1 - p) \cdot p'\bigr)
\;\ge\; p \cdot \hbin(1) + (1 - p) \cdot \hbin(p')
\;=\; (1 - p) \, \hbin(p'),
\end{align*}
where the last equality uses $\hbin(1) = 0$. This completes the proof.
\end{proof}
